\documentclass[11pt,a4paper]{article}

\usepackage{fontspec}
\usepackage{geometry}
\usepackage{hyperref}
\usepackage{longtable}
\usepackage{array}
\usepackage{booktabs}
\usepackage{enumitem}
\usepackage{xcolor}
\usepackage{fancyvrb}

\geometry{margin=25mm}
\setmainfont{Latin Modern Roman}
\setsansfont{Latin Modern Sans}
\setmonofont{Latin Modern Mono}
\hypersetup{
  colorlinks=true,
  linkcolor=blue!50!black,
  urlcolor=blue!50!black
}
\setlist[itemize]{itemsep=0.3em, topsep=0.4em}
\setlist[enumerate]{itemsep=0.3em, topsep=0.4em}

\DefineVerbatimEnvironment{codeblock}{Verbatim}{
  fontsize=\small,
  frame=single,
  rulecolor=\color{black!20}
}
\newcolumntype{L}[1]{>{\raggedright\arraybackslash}p{#1}}
\newcommand{\codeheading}[1]{\textnormal{\texttt{#1}}}

\title{svg2tex-rs Documentation}
\author{Eito Yoneyama}
\date{\today}

\begin{document}
\maketitle

\section{Overview}

\texttt{svg2tex-rs} provides the \texttt{svg2tex} command, which converts SVG files into either raw PDF literal operators or complete TeX source.
Its goal is to make SVG a practical input format for TeX workflows without giving up vector output
for common artwork.

The converter prefers native PDF constructs whenever possible. When an SVG subtree is better
represented by a raster intermediate image, the tool uses hybrid rendering instead of silently
dropping the feature. This keeps conversion behavior explicit and visually safer for real-world
documents.

\section{Release Scope}

Version 0.1.0 is suitable for public release as long as its scope is described honestly:

\begin{itemize}
\item It is strong on vector-preserving SVG to TeX conversion.
\item It supports a broad range of SVG features, including gradients, patterns, masks, clip paths,
filters, and embedded images.
\item It includes reproducibility controls for text and font resolution.
\item It may still use hybrid raster rendering for SVG subtrees that naturally require raster-domain
processing.
\end{itemize}

This means the right release message is not ``all SVG is always pure vector PDF'', but rather
``broad SVG support for TeX, with native vector output where practical and safe fallback where
necessary''.

\section{Installation}

\subsection{Install from crates.io with Cargo}

\begin{codeblock}
cargo install svg2tex-rs
\end{codeblock}

\subsection{Clone and Build Locally Without Installing}

\begin{codeblock}
git clone https://github.com/rice8y/svg2tex-rs.git
cd svg2tex-rs
just build
\end{codeblock}

The resulting binary is written to \texttt{target/release/svg2tex}.

\subsection{Clone and Install with just}

\begin{codeblock}
git clone https://github.com/rice8y/svg2tex-rs.git
cd svg2tex-rs
just install
\end{codeblock}

By default this installs the binary into \texttt{\$HOME/.local/bin}. To change the target:

\begin{codeblock}
LOCAL_BIN_DIR=/your/bin/path just install
\end{codeblock}

\section{Basic Usage}

\subsection{Emit Raw PDF Literal}

\begin{codeblock}
svg2tex --input drawing.svg --output drawing.literal
\end{codeblock}

This mode emits only the literal PDF operator stream.

\subsection{Emit a Standalone TeX Document}

\begin{codeblock}
svg2tex --input drawing.svg --tex --output drawing.tex
lualatex drawing.tex
\end{codeblock}

This is the recommended default when you want a directly compilable file.

\subsection{Emit a TeX Snippet}

\begin{codeblock}
svg2tex --input drawing.svg --tex --tex-format snippet --output drawing-snippet.tex
\end{codeblock}

This mode emits TeX definitions without a document wrapper, which is useful for inclusion into
another TeX source tree.

\section{CLI Reference}

\begin{longtable}{L{0.34\linewidth}L{0.60\linewidth}}
\toprule
\textbf{Option} & \textbf{Meaning} \\
\midrule
\endfirsthead
\multicolumn{2}{r}{\textit{CLI Reference continued from previous page}} \\
\toprule
\textbf{Option} & \textbf{Meaning} \\
\midrule
\endhead
\midrule
\multicolumn{2}{r}{\textit{Continued on next page}} \\
\endfoot
\bottomrule
\endlastfoot
\texttt{-i, --input <INPUT>} & Input SVG file path. \\
\texttt{-o, --output <OUTPUT>} & Output path. If omitted, output is written to standard output. \\
\texttt{--tex} & Emit TeX output instead of only PDF literal content. \\
\texttt{--tex-format <FORMAT>} & TeX wrapper format: \texttt{standalone}, \texttt{article}, or \texttt{snippet}. \\
\texttt{--embed-images} & Embed raster images referenced by the SVG. \\
\texttt{--no-system-fonts} & Disable host system font loading. \\
\texttt{--strict-fonts} & Error out if text conversion still depends on host-side font discovery. \\
\texttt{--report-fonts} & Print the font requirements observed in SVG text nodes. \\
\texttt{--font-family <NAME>} & Set the default font family used when the SVG omits one. \\
\texttt{--font-size <SIZE>} & Set the default font size used when the SVG omits one. \\
\texttt{--serif-family <NAME>} & Override the generic serif family. \\
\texttt{--sans-serif-family <NAME>} & Override the generic sans-serif family. \\
\texttt{--cursive-family <NAME>} & Override the generic cursive family. \\
\texttt{--fantasy-family <NAME>} & Override the generic fantasy family. \\
\texttt{--monospace-family <NAME>} & Override the generic monospace family. \\
\texttt{--font-file <PATH>} & Load an additional font file before parsing. May be repeated. \\
\texttt{--font-dir <PATH>} & Load all fonts from a directory before parsing. May be repeated. \\
\texttt{--strict} & Fail instead of hybrid-rendering unsupported SVG subtrees. \\
\texttt{--fallback-dpi <N>} & Set the raster fallback resolution. \\
\texttt{--engine <ENGINE>} & Choose \texttt{auto}, \texttt{pdftex}, \texttt{luatex}, \texttt{xetex}, \texttt{ptex}, or \texttt{uptex}. \\
\end{longtable}

\section{TeX Output Modes}

\subsection{\codeheading{standalone}}

This is the default mode when \texttt{--tex} is enabled. It emits a complete cropped document.
For direct compilation and quick inspection, this is the best default.

\subsection{\codeheading{article}}

This emits a complete article-class document. Use it when you want a less specialized wrapper
or when integrating into workflows that already assume the standard article layout.

\subsection{\codeheading{snippet}}

This mode emits only the TeX definitions needed to place the converted drawing. It is intended for
projects that already own the surrounding document structure.

\section{Rendering Strategy}

\subsection{Vector First}

The converter tries to preserve SVG semantics using native PDF constructs such as:

\begin{itemize}
\item path operators
\item clipping paths
\item shadings for gradients
\item patterns
\item form XObjects
\item soft masks
\item graphics states
\end{itemize}

\subsection{Hybrid Rendering}

Some SVG features are naturally raster-domain operations, especially certain filter pipelines.
In those situations, \texttt{svg2tex} rasterizes only the necessary subtree whenever possible.
This avoids both silent feature loss and unnecessary full-document rasterization.

\subsection{Strict Mode}

Use \texttt{--strict} when you want conversion to fail instead of using hybrid rendering:

\begin{codeblock}
svg2tex --input drawing.svg --tex --strict --output drawing.tex
\end{codeblock}

\section{Text and Font Reproducibility}

SVG text is flattened into paths during conversion. That makes the final TeX/PDF output self-contained
with respect to glyph outlines, but the outlines themselves still depend on which fonts were available
to the converter at parse time.

For reproducible builds:

\begin{enumerate}
\item Disable system font discovery with \texttt{--no-system-fonts}.
\item Load the required fonts explicitly with \texttt{--font-file} or \texttt{--font-dir}.
\item Use \texttt{--strict-fonts} to fail if a requested named font is missing.
\item Use \texttt{--report-fonts} while diagnosing text-heavy artwork.
\end{enumerate}

Recommended reproducible command:

\begin{codeblock}
svg2tex \
  --input poster.svg \
  --tex \
  --no-system-fonts \
  --strict-fonts \
  --font-file ./fonts/SourceSans3-Regular.ttf \
  --font-family "Source Sans 3" \
  --output poster.tex
\end{codeblock}

\section{Examples}

\subsection{Compile with LuaLaTeX}

\begin{codeblock}
svg2tex --input figure.svg --tex --engine luatex --output figure.tex
lualatex figure.tex
\end{codeblock}

\subsection{Compile with pdfTeX}

\begin{codeblock}
svg2tex --input figure.svg --tex --engine pdftex --output figure.tex
pdflatex figure.tex
\end{codeblock}

\subsection{DVI Workflow}

\begin{codeblock}
svg2tex --input figure.svg --tex --engine uptex --output figure.tex
uplatex figure.tex
dvipdfmx figure.dvi
\end{codeblock}

\section{Operational Notes}

\begin{itemize}
\item If exact reproduction of text is important, treat font configuration as part of the build input.
\item If exact vector-only output is important, use \texttt{--strict} and review failures explicitly.
\item If a document contains complex filters, hybrid rendering may be the correct and expected result.
\item For direct user-facing output, \texttt{--tex --tex-format standalone} is the most convenient mode.
\end{itemize}

\section{Conclusion}

\texttt{svg2tex-rs} v0.1.0 is a solid first public release for TeX users who want SVG support with a
serious bias toward native vector output. Its strengths are breadth of accepted SVG input, explicit
handling of difficult features, and practical controls for reproducible text conversion.

\end{document}
